\documentclass[10pt,a4paper]{amsart}
\usepackage{amsmath,amssymb,amsthm,graphicx}
% For commutative diagrams we use Xy-pic
\usepackage[all]{xy}

% We use 2cell for 2-commutative diagrams.
\xyoption{2cell}
% Theorem environments.
%
\theoremstyle{plain}
\newtheorem{theorem}[subsection]{Theorem}
\newtheorem{proposition}[subsection]{Proposition}
\newtheorem{lemma}[subsection]{Lemma}

\theoremstyle{definition}
\newtheorem{definition}[subsection]{Definition}
\newtheorem{example}[subsection]{Example}
\newtheorem{exercise}[subsection]{Exercise}
\newtheorem{situation}[subsection]{Situation}

\theoremstyle{remark}
\newtheorem{remark}[subsection]{Remark}
\newtheorem{remarks}[subsection]{Remarks}

\numberwithin{equation}{subsection}

% Macros
%
\def\lim{\mathop{\mathrm{lim}}\nolimits}
\def\colim{\mathop{\mathrm{colim}}\nolimits}
\def\Spec{\mathop{\mathrm{Spec}}}
\def\Hom{\mathop{\mathrm{Hom}}\nolimits}
\def\Ext{\mathop{\mathrm{Ext}}\nolimits}
\def\SheafHom{\mathop{\mathcal{H}\!\mathit{om}}\nolimits}
\def\SheafExt{\mathop{\mathcal{E}\!\mathit{xt}}\nolimits}
\def\Sch{\mathit{Sch}}
\def\Mor{\mathop{\mathrm{Mor}}\nolimits}
\def\Ob{\mathop{\mathrm{Ob}}\nolimits}
\def\Sh{\mathop{\mathit{Sh}}\nolimits}
\def\NL{\mathop{N\!L}\nolimits}
\def\CH{\mathop{\mathrm{CH}}\nolimits}
\def\proetale{{pro\text{-}\acute{e}tale}}
\def\etale{{\acute{e}tale}}
\def\QCoh{\mathit{QCoh}}
\def\Ker{\mathop{\mathrm{Ker}}}
\def\Im{\mathop{\mathrm{Im}}}
\def\Coker{\mathop{\mathrm{Coker}}}
\def\Coim{\mathop{\mathrm{Coim}}}

% Boxtimes
%
\DeclareMathSymbol{\boxtimes}{\mathbin}{AMSa}{"02}

%
% Macros for moduli stacks/spaces
%
\def\QCohstack{\mathcal{QC}\!\mathit{oh}}
\def\Cohstack{\mathcal{C}\!\mathit{oh}}
\def\Spacesstack{\mathcal{S}\!\mathit{paces}}
\def\Quotfunctor{\mathrm{Quot}}
\def\Hilbfunctor{\mathrm{Hilb}}
\def\Curvesstack{\mathcal{C}\!\mathit{urves}}
\def\Polarizedstack{\mathcal{P}\!\mathit{olarized}}
\def\Complexesstack{\mathcal{C}\!\mathit{omplexes}}
% \Pic is the operator that assigns to X its picard group, usage \Pic(X)
% \Picardstack_{X/B} denotes the Picard stack of X over B
% \Picardfunctor_{X/B} denotes the Picard functor of X over B
\def\Pic{\mathop{\mathrm{Pic}}\nolimits}
\def\Picardstack{\mathcal{P}\!\mathit{ic}}
\def\Picardfunctor{\mathrm{Pic}}
\def\Deformationcategory{\mathcal{D}\!\mathit{ef}}
\newenvironment{reference}{}{}
\newenvironment{slogan}{}{}
\newenvironment{history}{}{}
\usepackage{cleveref}
\crefname{originalsection}{original section}{original sections}
\Crefname{originalsection}{Original section}{Original sections}
\makeatletter
\newlabel{spaces-more-morphisms-section-relatively-perfect}{{0DKM}{source}{}{sourcepointer-1}{}}
\newlabel{spaces-more-morphisms-section-relatively-perfect@cref}{{[originalsection][0][]0DKM}{[1][0][]source}{}{}{}}
\makeatother
\crefname{originaldefinition}{original definition}{original definitions}
\Crefname{originaldefinition}{Original definition}{Original definitions}
\makeatletter
\newlabel{spaces-more-morphisms-definition-relatively-perfect}{{0DKN}{source}{}{sourcepointer-2}{}}
\newlabel{spaces-more-morphisms-definition-relatively-perfect@cref}{{[originaldefinition][0][]0DKN}{[1][0][]source}{}{}{}}
\makeatother
\crefname{originallemma}{original lemma}{original lemmas}
\Crefname{originallemma}{Original lemma}{Original lemmas}
\makeatletter
\newlabel{sites-cohomology-lemma-pseudo-coherent-pullback}{{08H4}{source}{}{sourcepointer-3}{}}
\newlabel{sites-cohomology-lemma-pseudo-coherent-pullback@cref}{{[originallemma][0][]08H4}{[1][0][]source}{}{}{}}
\makeatother
\crefname{originallemma}{original lemma}{original lemmas}
\Crefname{originallemma}{Original lemma}{Original lemmas}
\makeatletter
\newlabel{spaces-more-morphisms-lemma-base-change-relatively-perfect}{{0DKT}{source}{}{sourcepointer-4}{}}
\newlabel{spaces-more-morphisms-lemma-base-change-relatively-perfect@cref}{{[originallemma][0][]0DKT}{[1][0][]source}{}{}{}}
\makeatother
\crefname{originallemma}{original lemma}{original lemmas}
\Crefname{originallemma}{Original lemma}{Original lemmas}
\makeatletter
\newlabel{perfect-lemma-affine-morphism-and-hom-out-of-perfect}{{0AA7}{source}{}{sourcepointer-5}{}}
\newlabel{perfect-lemma-affine-morphism-and-hom-out-of-perfect@cref}{{[originallemma][0][]0AA7}{[1][0][]source}{}{}{}}
\makeatother
\crefname{originallemma}{original lemma}{original lemmas}
\Crefname{originallemma}{Original lemma}{Original lemmas}
\makeatletter
\newlabel{spaces-more-morphisms-lemma-compute-ext-rel-perfect}{{0DKY}{source}{}{sourcepointer-6}{}}
\newlabel{spaces-more-morphisms-lemma-compute-ext-rel-perfect@cref}{{[originallemma][0][]0DKY}{[1][0][]source}{}{}{}}
\makeatother
\crefname{originallemma}{original lemma}{original lemmas}
\Crefname{originallemma}{Original lemma}{Original lemmas}
\makeatletter
\newlabel{perfect-lemma-jump-loci}{{0BDI}{source}{}{sourcepointer-7}{}}
\newlabel{perfect-lemma-jump-loci@cref}{{[originallemma][0][]0BDI}{[1][0][]source}{}{}{}}
\makeatother
\crefname{originallemma}{original lemma}{original lemmas}
\Crefname{originallemma}{Original lemma}{Original lemmas}
\makeatletter
\newlabel{sites-cohomology-lemma-section-RHom-over-U}{{08JA}{source}{}{sourcepointer-8}{}}
\newlabel{sites-cohomology-lemma-section-RHom-over-U@cref}{{[originallemma][0][]08JA}{[1][0][]source}{}{}{}}
\makeatother
\crefname{originallemma}{original lemma}{original lemmas}
\Crefname{originallemma}{Original lemma}{Original lemmas}
\makeatletter
\newlabel{sites-cohomology-lemma-Leray-unbounded}{{0D6H}{source}{}{sourcepointer-9}{}}
\newlabel{sites-cohomology-lemma-Leray-unbounded@cref}{{[originallemma][0][]0D6H}{[1][0][]source}{}{}{}}
\makeatother
\crefname{originallemma}{original lemma}{original lemmas}
\Crefname{originallemma}{Original lemma}{Original lemmas}
\makeatletter
\newlabel{derived-lemma-negative-vanishing}{{05TC}{source}{}{sourcepointer-10}{}}
\newlabel{derived-lemma-negative-vanishing@cref}{{[originallemma][0][]05TC}{[1][0][]source}{}{}{}}
\makeatother
\crefname{originallemma}{original lemma}{original lemmas}
\Crefname{originallemma}{Original lemma}{Original lemmas}
\makeatletter
\newlabel{more-algebra-lemma-lift-pseudo-coherent-from-residue-field}{{0BCC}{source}{}{sourcepointer-11}{}}
\newlabel{more-algebra-lemma-lift-pseudo-coherent-from-residue-field@cref}{{[originallemma][0][]0BCC}{[1][0][]source}{}{}{}}
\makeatother
\crefname{originallemma}{original lemma}{original lemmas}
\Crefname{originallemma}{Original lemma}{Original lemmas}
\makeatletter
\newlabel{spaces-simplicial-lemma-fppf-neg-ext-zero-hom}{{0DL9}{source}{}{sourcepointer-12}{}}
\newlabel{spaces-simplicial-lemma-fppf-neg-ext-zero-hom@cref}{{[originallemma][0][]0DL9}{[1][0][]source}{}{}{}}
\makeatother
\crefname{originallemma}{original lemma}{original lemmas}
\Crefname{originallemma}{Original lemma}{Original lemmas}
\makeatletter
\newlabel{bootstrap-lemma-locally-algebraic-space-finite-type}{{04U0}{source}{}{sourcepointer-13}{}}
\newlabel{bootstrap-lemma-locally-algebraic-space-finite-type@cref}{{[originallemma][0][]04U0}{[1][0][]source}{}{}{}}
\makeatother
\crefname{originallemma}{original lemma}{original lemmas}
\Crefname{originallemma}{Original lemma}{Original lemmas}
\makeatletter
\newlabel{spaces-morphisms-lemma-affine-local}{{03WG}{source}{}{sourcepointer-14}{}}
\newlabel{spaces-morphisms-lemma-affine-local@cref}{{[originallemma][0][]03WG}{[1][0][]source}{}{}{}}
\makeatother
\crefname{originallemma}{original lemma}{original lemmas}
\Crefname{originallemma}{Original lemma}{Original lemmas}
\makeatletter
\newlabel{spaces-morphisms-lemma-finite-presentation-local}{{0410}{source}{}{sourcepointer-15}{}}
\newlabel{spaces-morphisms-lemma-finite-presentation-local@cref}{{[originallemma][0][]0410}{[1][0][]source}{}{}{}}
\makeatother
\crefname{originallemma}{original lemma}{original lemmas}
\Crefname{originallemma}{Original lemma}{Original lemmas}
\makeatletter
\newlabel{more-algebra-lemma-pseudo-coherent}{{064U}{source}{}{sourcepointer-16}{}}
\newlabel{more-algebra-lemma-pseudo-coherent@cref}{{[originallemma][0][]064U}{[1][0][]source}{}{}{}}
\makeatother
\crefname{originaldefinition}{original definition}{original definitions}
\Crefname{originaldefinition}{Original definition}{Original definitions}
\makeatletter
\newlabel{categories-definition-fibred-groupoids}{{003T}{source}{}{sourcepointer-17}{}}
\newlabel{categories-definition-fibred-groupoids@cref}{{[originaldefinition][0][]003T}{[1][0][]source}{}{}{}}
\makeatother
\crefname{originallemma}{original lemma}{original lemmas}
\Crefname{originallemma}{Original lemma}{Original lemmas}
\makeatletter
\newlabel{algebraic-lemma-representable-diagonal}{{045G}{source}{}{sourcepointer-18}{}}
\newlabel{algebraic-lemma-representable-diagonal@cref}{{[originallemma][0][]045G}{[1][0][]source}{}{}{}}
\makeatother
\crefname{originalproposition}{original proposition}{original propositions}
\Crefname{originalproposition}{Original proposition}{Original propositions}
\makeatletter
\newlabel{proposition-isom}{{08K9}{source}{}{sourcepointer-19}{}}
\newlabel{proposition-isom@cref}{{[originalproposition][0][]08K9}{[1][0][]source}{}{}{}}
\makeatother
\crefname{originallemma}{original lemma}{original lemmas}
\Crefname{originallemma}{Original lemma}{Original lemmas}
\makeatletter
\newlabel{spaces-descent-lemma-fpqc-universal-effective-epimorphisms}{{04P2}{source}{}{sourcepointer-20}{}}
\newlabel{spaces-descent-lemma-fpqc-universal-effective-epimorphisms@cref}{{[originallemma][0][]04P2}{[1][0][]source}{}{}{}}
\makeatother
\crefname{originallemma}{original lemma}{original lemmas}
\Crefname{originallemma}{Original lemma}{Original lemmas}
\makeatletter
\newlabel{spaces-simplicial-lemma-fppf-glue-neg-ext-zero}{{0DLA}{source}{}{sourcepointer-21}{}}
\newlabel{spaces-simplicial-lemma-fppf-glue-neg-ext-zero@cref}{{[originallemma][0][]0DLA}{[1][0][]source}{}{}{}}
\makeatother
\crefname{originallemma}{original lemma}{original lemmas}
\Crefname{originallemma}{Original lemma}{Original lemmas}
\makeatletter
\newlabel{spaces-more-morphisms-lemma-pseudo-coherent-descends-fpqc}{{0DL1}{source}{}{sourcepointer-22}{}}
\newlabel{spaces-more-morphisms-lemma-pseudo-coherent-descends-fpqc@cref}{{[originallemma][0][]0DL1}{[1][0][]source}{}{}{}}
\makeatother
\crefname{originallemma}{original lemma}{original lemmas}
\Crefname{originallemma}{Original lemma}{Original lemmas}
\makeatletter
\newlabel{spaces-more-morphisms-lemma-tor-amplitude-descends-fppf}{{0DL2}{source}{}{sourcepointer-23}{}}
\newlabel{spaces-more-morphisms-lemma-tor-amplitude-descends-fppf@cref}{{[originallemma][0][]0DL2}{[1][0][]source}{}{}{}}
\makeatother
\crefname{originalsection}{original section}{original sections}
\Crefname{originalsection}{Original section}{Original sections}
\makeatletter
\newlabel{algebraic-section-split}{{04SU}{source}{}{sourcepointer-24}{}}
\newlabel{algebraic-section-split@cref}{{[originalsection][0][]04SU}{[1][0][]source}{}{}{}}
\makeatother
\crefname{originalsection}{original section}{original sections}
\Crefname{originalsection}{Original section}{Original sections}
\makeatletter
\newlabel{examples-stacks-section-stack-associated-to-sheaf}{{0305}{source}{}{sourcepointer-25}{}}
\newlabel{examples-stacks-section-stack-associated-to-sheaf@cref}{{[originalsection][0][]0305}{[1][0][]source}{}{}{}}
\makeatother
\crefname{originaldefinition}{original definition}{original definitions}
\Crefname{originaldefinition}{Original definition}{Original definitions}
\makeatletter
\newlabel{artin-definition-limit-preserving}{{07XL}{source}{}{sourcepointer-26}{}}
\newlabel{artin-definition-limit-preserving@cref}{{[originaldefinition][0][]07XL}{[1][0][]source}{}{}{}}
\makeatother
\crefname{originalsection}{original section}{original sections}
\Crefname{originalsection}{Original section}{Original sections}
\makeatletter
\newlabel{spaces-limits-remark-limit-preserving}{{05N0}{source}{}{sourcepointer-27}{}}
\newlabel{spaces-limits-remark-limit-preserving@cref}{{[originalsection][0][]05N0}{[1][0][]source}{}{}{}}
\makeatother
\crefname{originallemma}{original lemma}{original lemmas}
\Crefname{originallemma}{Original lemma}{Original lemmas}
\makeatletter
\newlabel{spaces-more-morphisms-lemma-descend-relatively-perfect}{{0DKW}{source}{}{sourcepointer-28}{}}
\newlabel{spaces-more-morphisms-lemma-descend-relatively-perfect@cref}{{[originallemma][0][]0DKW}{[1][0][]source}{}{}{}}
\makeatother
\crefname{originallemma}{original lemma}{original lemmas}
\Crefname{originallemma}{Original lemma}{Original lemmas}
\makeatletter
\newlabel{limits-lemma-limit-contained-in-constructible}{{05F4}{source}{}{sourcepointer-29}{}}
\newlabel{limits-lemma-limit-contained-in-constructible@cref}{{[originallemma][0][]05F4}{[1][0][]source}{}{}{}}
\makeatother
\crefname{originallemma}{original lemma}{original lemmas}
\Crefname{originallemma}{Original lemma}{Original lemmas}
\makeatletter
\newlabel{more-morphisms-lemma-pushout-along-thickening}{{07RT}{source}{}{sourcepointer-30}{}}
\newlabel{more-morphisms-lemma-pushout-along-thickening@cref}{{[originallemma][0][]07RT}{[1][0][]source}{}{}{}}
\makeatother
\crefname{originallemma}{original lemma}{original lemmas}
\Crefname{originallemma}{Original lemma}{Original lemmas}
\makeatletter
\newlabel{spaces-pushouts-lemma-pushout-along-thickening-schemes}{{07SY}{source}{}{sourcepointer-31}{}}
\newlabel{spaces-pushouts-lemma-pushout-along-thickening-schemes@cref}{{[originallemma][0][]07SY}{[1][0][]source}{}{}{}}
\makeatother
\crefname{originallemma}{original lemma}{original lemmas}
\Crefname{originallemma}{Original lemma}{Original lemmas}
\makeatletter
\newlabel{artin-lemma-pushout}{{07WN}{source}{}{sourcepointer-32}{}}
\newlabel{artin-lemma-pushout@cref}{{[originallemma][0][]07WN}{[1][0][]source}{}{}{}}
\makeatother
\crefname{originallemma}{original lemma}{original lemmas}
\Crefname{originallemma}{Original lemma}{Original lemmas}
\makeatletter
\newlabel{spaces-pushouts-lemma-equivalence-categories-spaces-pushout-flat}{{07W3}{source}{}{sourcepointer-33}{}}
\newlabel{spaces-pushouts-lemma-equivalence-categories-spaces-pushout-flat@cref}{{[originallemma][0][]07W3}{[1][0][]source}{}{}{}}
\makeatother
\crefname{originallemma}{original lemma}{original lemmas}
\Crefname{originallemma}{Original lemma}{Original lemmas}
\makeatletter
\newlabel{spaces-pushouts-lemma-pushout-along-thickening-derived}{{0DL7}{source}{}{sourcepointer-34}{}}
\newlabel{spaces-pushouts-lemma-pushout-along-thickening-derived@cref}{{[originallemma][0][]0DL7}{[1][0][]source}{}{}{}}
\makeatother
\crefname{originallemma}{original lemma}{original lemmas}
\Crefname{originallemma}{Original lemma}{Original lemmas}
\makeatletter
\newlabel{spaces-more-morphisms-lemma-thickening-relatively-perfect}{{0DL4}{source}{}{sourcepointer-35}{}}
\newlabel{spaces-more-morphisms-lemma-thickening-relatively-perfect@cref}{{[originallemma][0][]0DL4}{[1][0][]source}{}{}{}}
\makeatother
\crefname{originalsection}{original section}{original sections}
\Crefname{originalsection}{Original section}{Original sections}
\makeatletter
\newlabel{artin-section-tangent-spaces}{{07WY}{source}{}{sourcepointer-36}{}}
\newlabel{artin-section-tangent-spaces@cref}{{[originalsection][0][]07WY}{[1][0][]source}{}{}{}}
\makeatother
\crefname{originalsection}{original section}{original sections}
\Crefname{originalsection}{Original section}{Original sections}
\makeatletter
\newlabel{artin-section-RS-star}{{0CXN}{source}{}{sourcepointer-37}{}}
\newlabel{artin-section-RS-star@cref}{{[originalsection][0][]0CXN}{[1][0][]source}{}{}{}}
\makeatother
\crefname{originallemma}{original lemma}{original lemmas}
\Crefname{originallemma}{Original lemma}{Original lemmas}
\makeatletter
\newlabel{artin-lemma-deformation-category}{{07WU}{source}{}{sourcepointer-38}{}}
\newlabel{artin-lemma-deformation-category@cref}{{[originallemma][0][]07WU}{[1][0][]source}{}{}{}}
\makeatother
\crefname{originallemma}{original lemma}{original lemmas}
\Crefname{originallemma}{Original lemma}{Original lemmas}
\makeatletter
\newlabel{artin-lemma-fibre-product-tangent-spaces}{{07X2}{source}{}{sourcepointer-39}{}}
\newlabel{artin-lemma-fibre-product-tangent-spaces@cref}{{[originallemma][0][]07X2}{[1][0][]source}{}{}{}}
\makeatother
\crefname{originallemma}{original lemma}{original lemmas}
\Crefname{originallemma}{Original lemma}{Original lemmas}
\makeatletter
\newlabel{artin-lemma-finite-dimension}{{07X1}{source}{}{sourcepointer-40}{}}
\newlabel{artin-lemma-finite-dimension@cref}{{[originallemma][0][]07X1}{[1][0][]source}{}{}{}}
\makeatother
\crefname{originalsection}{original section}{original sections}
\Crefname{originalsection}{Original section}{Original sections}
\makeatletter
\newlabel{spaces-more-morphisms-section-thickenings}{{05ZJ}{source}{}{sourcepointer-41}{}}
\newlabel{spaces-more-morphisms-section-thickenings@cref}{{[originalsection][0][]05ZJ}{[1][0][]source}{}{}{}}
\makeatother
\crefname{originallemma}{original lemma}{original lemmas}
\Crefname{originallemma}{Original lemma}{Original lemmas}
\makeatletter
\newlabel{defos-lemma-first-order-defos-complex}{{0DIZ}{source}{}{sourcepointer-42}{}}
\newlabel{defos-lemma-first-order-defos-complex@cref}{{[originallemma][0][]0DIZ}{[1][0][]source}{}{}{}}
\makeatother
\crefname{originallemma}{original lemma}{original lemmas}
\Crefname{originallemma}{Original lemma}{Original lemmas}
\makeatletter
\newlabel{spaces-perfect-lemma-identify-pseudo-coherent-noetherian}{{08IK}{source}{}{sourcepointer-43}{}}
\newlabel{spaces-perfect-lemma-identify-pseudo-coherent-noetherian@cref}{{[originallemma][0][]08IK}{[1][0][]source}{}{}{}}
\makeatother
\crefname{originallemma}{original lemma}{original lemmas}
\Crefname{originallemma}{Original lemma}{Original lemmas}
\makeatletter
\newlabel{spaces-perfect-lemma-ext-finite}{{0D0T}{source}{}{sourcepointer-44}{}}
\newlabel{spaces-perfect-lemma-ext-finite@cref}{{[originallemma][0][]0D0T}{[1][0][]source}{}{}{}}
\makeatother
\crefname{originalsection}{original section}{original sections}
\Crefname{originalsection}{Original section}{Original sections}
\makeatletter
\newlabel{artin-remark-strong-effectiveness}{{0CXT}{source}{}{sourcepointer-45}{}}
\newlabel{artin-remark-strong-effectiveness@cref}{{[originalsection][0][]0CXT}{[1][0][]source}{}{}{}}
\makeatother
\crefname{originaltheorem}{original theorem}{original theorems}
\Crefname{originaltheorem}{Original theorem}{Original theorems}
\makeatletter
\newlabel{spaces-flat-theorem-existence-derived}{{0DIQ}{source}{}{sourcepointer-46}{}}
\newlabel{spaces-flat-theorem-existence-derived@cref}{{[originaltheorem][0][]0DIQ}{[1][0][]source}{}{}{}}
\makeatother
\crefname{originallemma}{original lemma}{original lemmas}
\Crefname{originallemma}{Original lemma}{Original lemmas}
\makeatletter
\newlabel{more-algebra-lemma-limit-henselian}{{0CT7}{source}{}{sourcepointer-47}{}}
\newlabel{more-algebra-lemma-limit-henselian@cref}{{[originallemma][0][]0CT7}{[1][0][]source}{}{}{}}
\makeatother
\crefname{originallemma}{original lemma}{original lemmas}
\Crefname{originallemma}{Original lemma}{Original lemmas}
\makeatletter
\newlabel{spaces-more-morphisms-lemma-henselian-relatively-perfect}{{0DL5}{source}{}{sourcepointer-48}{}}
\newlabel{spaces-more-morphisms-lemma-henselian-relatively-perfect@cref}{{[originallemma][0][]0DL5}{[1][0][]source}{}{}{}}
\makeatother
\crefname{originalsection}{original section}{original sections}
\Crefname{originalsection}{Original section}{Original sections}
\makeatletter
\newlabel{algebraic-section-representable-properties}{{03YJ}{source}{}{sourcepointer-49}{}}
\newlabel{algebraic-section-representable-properties@cref}{{[originalsection][0][]03YJ}{[1][0][]source}{}{}{}}
\makeatother
\crefname{originalsection}{original section}{original sections}
\Crefname{originalsection}{Original section}{Original sections}
\makeatletter
\newlabel{algebraic-section-change-base-scheme}{{04X4}{source}{}{sourcepointer-50}{}}
\newlabel{algebraic-section-change-base-scheme@cref}{{[originalsection][0][]04X4}{[1][0][]source}{}{}{}}
\makeatother
\crefname{originallemma}{original lemma}{original lemmas}
\Crefname{originallemma}{Original lemma}{Original lemmas}
\makeatletter
\newlabel{spaces-limits-lemma-descend-finite-presentation}{{07SK}{source}{}{sourcepointer-51}{}}
\newlabel{spaces-limits-lemma-descend-finite-presentation@cref}{{[originallemma][0][]07SK}{[1][0][]source}{}{}{}}
\makeatother
\crefname{originallemma}{original lemma}{original lemmas}
\Crefname{originallemma}{Original lemma}{Original lemmas}
\makeatletter
\newlabel{spaces-limits-lemma-descend-flat}{{08K0}{source}{}{sourcepointer-52}{}}
\newlabel{spaces-limits-lemma-descend-flat@cref}{{[originallemma][0][]08K0}{[1][0][]source}{}{}{}}
\makeatother
\crefname{originallemma}{original lemma}{original lemmas}
\Crefname{originallemma}{Original lemma}{Original lemmas}
\makeatletter
\newlabel{spaces-limits-lemma-eventually-proper}{{08K1}{source}{}{sourcepointer-53}{}}
\newlabel{spaces-limits-lemma-eventually-proper@cref}{{[originallemma][0][]08K1}{[1][0][]source}{}{}{}}
\makeatother
\crefname{originallemma}{original lemma}{original lemmas}
\Crefname{originallemma}{Original lemma}{Original lemmas}
\makeatletter
\newlabel{artin-lemma-diagonal-representable}{{07Y4}{source}{}{sourcepointer-54}{}}
\newlabel{artin-lemma-diagonal-representable@cref}{{[originallemma][0][]07Y4}{[1][0][]source}{}{}{}}
\makeatother
\crefname{originalproposition}{original proposition}{original propositions}
\Crefname{originalproposition}{Original proposition}{Original propositions}
\makeatletter
\newlabel{more-algebra-proposition-ubiquity-G-ring}{{07PX}{source}{}{sourcepointer-55}{}}
\newlabel{more-algebra-proposition-ubiquity-G-ring@cref}{{[originalproposition][0][]07PX}{[1][0][]source}{}{}{}}
\makeatother
\crefname{originalsection}{original section}{original sections}
\Crefname{originalsection}{Original section}{Original sections}
\makeatletter
\newlabel{artin-section-axioms}{{07XJ}{source}{}{sourcepointer-56}{}}
\newlabel{artin-section-axioms@cref}{{[originalsection][0][]07XJ}{[1][0][]source}{}{}{}}
\makeatother
\crefname{originallemma}{original lemma}{original lemmas}
\Crefname{originallemma}{Original lemma}{Original lemmas}
\makeatletter
\newlabel{artin-lemma-SGE-implies-openness-versality}{{0CXU}{source}{}{sourcepointer-57}{}}
\newlabel{artin-lemma-SGE-implies-openness-versality@cref}{{[originallemma][0][]0CXU}{[1][0][]source}{}{}{}}
\makeatother
\title{Algebraicity of the moduli stack of relatively perfect universally gluable complexes}
\author{The Stacks Project Authors}
\date{stacks-project:quot.tex:local-snapshot-20261011; GNU FDL 1.2 or later}
\begin{document}
\maketitle
\noindent\textit{Editorial note.} Only the source-defined universally gluable relatively perfect objects are parameterized; this is not a moduli stack of all arbitrary complexes.\par\medskip
\noindent\textit{Editorial note.} The sheaf-moduli special cases are not counted as additional tasks from the same chain.\par\medskip
\noindent\textit{Allowed standard prerequisites (editorial).}\begin{itemize}
\item Explicit prerequisite boundary. Derived categories of algebraic spaces, relative perfection and negativeExt vanishing, formal deformation theory and Artin representability. The local stack/diagonal/limits/Rim–Schlessinger/tangent/effectiveness chain is retained author proof rather than granted as the main conclusion.
\end{itemize}\medskip
\ifdefined\hypertarget\hypertarget{sourcepointer-1}{}\fi
\noindent\textit{Source reference.} Original section 0DKM: focused-seeds/AGP/sources/stacks-snapshot-20261011/tags.txt, source lines 17017--17017. Stable Stacks tag for an expressly allowed foundational or notational input; the counted construction remains the two concrete core obligations. The retained local core is not replaced by this pointer. See standard\_prerequisites.\par\medskip
\ifdefined\hypertarget\hypertarget{sourcepointer-2}{}\fi
\noindent\textit{Source reference.} Original definition 0DKN: focused-seeds/AGP/sources/stacks-snapshot-20261011/tags.txt, source lines 17018--17018. Stable Stacks tag for an expressly allowed foundational or notational input; the counted construction remains the two concrete core obligations. The retained local core is not replaced by this pointer. See standard\_prerequisites.\par\medskip
\ifdefined\hypertarget\hypertarget{sourcepointer-3}{}\fi
\noindent\textit{Source reference.} Original lemma 08H4: focused-seeds/AGP/sources/stacks-snapshot-20261011/tags.txt, source lines 10767--10767. Stable Stacks tag for an expressly allowed foundational or notational input; the counted construction remains the two concrete core obligations. The retained local core is not replaced by this pointer. See standard\_prerequisites.\par\medskip
\ifdefined\hypertarget\hypertarget{sourcepointer-4}{}\fi
\noindent\textit{Source reference.} Original lemma 0DKT: focused-seeds/AGP/sources/stacks-snapshot-20261011/tags.txt, source lines 17023--17023. Stable Stacks tag for an expressly allowed foundational or notational input; the counted construction remains the two concrete core obligations. The retained local core is not replaced by this pointer. See standard\_prerequisites.\par\medskip
\ifdefined\hypertarget\hypertarget{sourcepointer-5}{}\fi
\noindent\textit{Source reference.} Original lemma 0AA7: focused-seeds/AGP/sources/stacks-snapshot-20261011/tags.txt, source lines 12975--12975. Stable Stacks tag for an expressly allowed foundational or notational input; the counted construction remains the two concrete core obligations. The retained local core is not replaced by this pointer. See standard\_prerequisites.\par\medskip
\ifdefined\hypertarget\hypertarget{sourcepointer-6}{}\fi
\noindent\textit{Source reference.} Original lemma 0DKY: focused-seeds/AGP/sources/stacks-snapshot-20261011/tags.txt, source lines 17028--17028. Stable Stacks tag for an expressly allowed foundational or notational input; the counted construction remains the two concrete core obligations. The retained local core is not replaced by this pointer. See standard\_prerequisites.\par\medskip
\ifdefined\hypertarget\hypertarget{sourcepointer-7}{}\fi
\noindent\textit{Source reference.} Original lemma 0BDI: focused-seeds/AGP/sources/stacks-snapshot-20261011/tags.txt, source lines 14316--14316. Stable Stacks tag for an expressly allowed foundational or notational input; the counted construction remains the two concrete core obligations. The retained local core is not replaced by this pointer. See standard\_prerequisites.\par\medskip
\ifdefined\hypertarget\hypertarget{sourcepointer-8}{}\fi
\noindent\textit{Source reference.} Original lemma 08JA: focused-seeds/AGP/sources/stacks-snapshot-20261011/tags.txt, source lines 10843--10843. Stable Stacks tag for an expressly allowed foundational or notational input; the counted construction remains the two concrete core obligations. The retained local core is not replaced by this pointer. See standard\_prerequisites.\par\medskip
\ifdefined\hypertarget\hypertarget{sourcepointer-9}{}\fi
\noindent\textit{Source reference.} Original lemma 0D6H: focused-seeds/AGP/sources/stacks-snapshot-20261011/tags.txt, source lines 16522--16522. Stable Stacks tag for an expressly allowed foundational or notational input; the counted construction remains the two concrete core obligations. The retained local core is not replaced by this pointer. See standard\_prerequisites.\par\medskip
\ifdefined\hypertarget\hypertarget{sourcepointer-10}{}\fi
\noindent\textit{Source reference.} Original lemma 05TC: focused-seeds/AGP/sources/stacks-snapshot-20261011/tags.txt, source lines 7476--7476. Stable Stacks tag for an expressly allowed foundational or notational input; the counted construction remains the two concrete core obligations. The retained local core is not replaced by this pointer. See standard\_prerequisites.\par\medskip
\ifdefined\hypertarget\hypertarget{sourcepointer-11}{}\fi
\noindent\textit{Source reference.} Original lemma 0BCC: focused-seeds/AGP/sources/stacks-snapshot-20261011/tags.txt, source lines 14275--14275. Stable Stacks tag for an expressly allowed foundational or notational input; the counted construction remains the two concrete core obligations. The retained local core is not replaced by this pointer. See standard\_prerequisites.\par\medskip
\ifdefined\hypertarget\hypertarget{sourcepointer-12}{}\fi
\noindent\textit{Source reference.} Original lemma 0DL9: focused-seeds/AGP/sources/stacks-snapshot-20261011/tags.txt, source lines 17039--17039. Stable Stacks tag for an expressly allowed foundational or notational input; the counted construction remains the two concrete core obligations. The retained local core is not replaced by this pointer. See standard\_prerequisites.\par\medskip
\ifdefined\hypertarget\hypertarget{sourcepointer-13}{}\fi
\noindent\textit{Source reference.} Original lemma 04U0: focused-seeds/AGP/sources/stacks-snapshot-20261011/tags.txt, source lines 6273--6273. Stable Stacks tag for an expressly allowed foundational or notational input; the counted construction remains the two concrete core obligations. The retained local core is not replaced by this pointer. See standard\_prerequisites.\par\medskip
\ifdefined\hypertarget\hypertarget{sourcepointer-14}{}\fi
\noindent\textit{Source reference.} Original lemma 03WG: focused-seeds/AGP/sources/stacks-snapshot-20261011/tags.txt, source lines 5111--5111. Stable Stacks tag for an expressly allowed foundational or notational input; the counted construction remains the two concrete core obligations. The retained local core is not replaced by this pointer. See standard\_prerequisites.\par\medskip
\ifdefined\hypertarget\hypertarget{sourcepointer-15}{}\fi
\noindent\textit{Source reference.} Original lemma 0410: focused-seeds/AGP/sources/stacks-snapshot-20261011/tags.txt, source lines 5275--5275. Stable Stacks tag for an expressly allowed foundational or notational input; the counted construction remains the two concrete core obligations. The retained local core is not replaced by this pointer. See standard\_prerequisites.\par\medskip
\ifdefined\hypertarget\hypertarget{sourcepointer-16}{}\fi
\noindent\textit{Source reference.} Original lemma 064U: focused-seeds/AGP/sources/stacks-snapshot-20261011/tags.txt, source lines 7878--7878. Stable Stacks tag for an expressly allowed foundational or notational input; the counted construction remains the two concrete core obligations. The retained local core is not replaced by this pointer. See standard\_prerequisites.\par\medskip
\ifdefined\hypertarget\hypertarget{sourcepointer-17}{}\fi
\noindent\textit{Source reference.} Original definition 003T: focused-seeds/AGP/sources/stacks-snapshot-20261011/tags.txt, source lines 147--147. Stable Stacks tag for an expressly allowed foundational or notational input; the counted construction remains the two concrete core obligations. The retained local core is not replaced by this pointer. See standard\_prerequisites.\par\medskip
\ifdefined\hypertarget\hypertarget{sourcepointer-18}{}\fi
\noindent\textit{Source reference.} Original lemma 045G: focused-seeds/AGP/sources/stacks-snapshot-20261011/tags.txt, source lines 5435--5435. Stable Stacks tag for an expressly allowed foundational or notational input; the counted construction remains the two concrete core obligations. The retained local core is not replaced by this pointer. See standard\_prerequisites.\par\medskip
\ifdefined\hypertarget\hypertarget{sourcepointer-19}{}\fi
\noindent\textit{Source reference.} Original proposition 08K9: focused-seeds/AGP/sources/stacks-snapshot-20261011/tags.txt, source lines 10877--10877. Stable Stacks tag for an expressly allowed foundational or notational input; the counted construction remains the two concrete core obligations. The retained local core is not replaced by this pointer. See standard\_prerequisites.\par\medskip
\ifdefined\hypertarget\hypertarget{sourcepointer-20}{}\fi
\noindent\textit{Source reference.} Original lemma 04P2: focused-seeds/AGP/sources/stacks-snapshot-20261011/tags.txt, source lines 6100--6100. Stable Stacks tag for an expressly allowed foundational or notational input; the counted construction remains the two concrete core obligations. The retained local core is not replaced by this pointer. See standard\_prerequisites.\par\medskip
\ifdefined\hypertarget\hypertarget{sourcepointer-21}{}\fi
\noindent\textit{Source reference.} Original lemma 0DLA: focused-seeds/AGP/sources/stacks-snapshot-20261011/tags.txt, source lines 17040--17040. Stable Stacks tag for an expressly allowed foundational or notational input; the counted construction remains the two concrete core obligations. The retained local core is not replaced by this pointer. See standard\_prerequisites.\par\medskip
\ifdefined\hypertarget\hypertarget{sourcepointer-22}{}\fi
\noindent\textit{Source reference.} Original lemma 0DL1: focused-seeds/AGP/sources/stacks-snapshot-20261011/tags.txt, source lines 17031--17031. Stable Stacks tag for an expressly allowed foundational or notational input; the counted construction remains the two concrete core obligations. The retained local core is not replaced by this pointer. See standard\_prerequisites.\par\medskip
\ifdefined\hypertarget\hypertarget{sourcepointer-23}{}\fi
\noindent\textit{Source reference.} Original lemma 0DL2: focused-seeds/AGP/sources/stacks-snapshot-20261011/tags.txt, source lines 17032--17032. Stable Stacks tag for an expressly allowed foundational or notational input; the counted construction remains the two concrete core obligations. The retained local core is not replaced by this pointer. See standard\_prerequisites.\par\medskip
\ifdefined\hypertarget\hypertarget{sourcepointer-24}{}\fi
\noindent\textit{Source reference.} Original section 04SU: focused-seeds/AGP/sources/stacks-snapshot-20261011/tags.txt, source lines 6232--6232. Stable Stacks tag for an expressly allowed foundational or notational input; the counted construction remains the two concrete core obligations. The retained local core is not replaced by this pointer. See standard\_prerequisites.\par\medskip
\ifdefined\hypertarget\hypertarget{sourcepointer-25}{}\fi
\noindent\textit{Source reference.} Original section 0305: focused-seeds/AGP/sources/stacks-snapshot-20261011/tags.txt, source lines 3936--3936. Stable Stacks tag for an expressly allowed foundational or notational input; the counted construction remains the two concrete core obligations. The retained local core is not replaced by this pointer. See standard\_prerequisites.\par\medskip
\ifdefined\hypertarget\hypertarget{sourcepointer-26}{}\fi
\noindent\textit{Source reference.} Original definition 07XL: focused-seeds/AGP/sources/stacks-snapshot-20261011/tags.txt, source lines 10081--10081. Stable Stacks tag for an expressly allowed foundational or notational input; the counted construction remains the two concrete core obligations. The retained local core is not replaced by this pointer. See standard\_prerequisites.\par\medskip
\ifdefined\hypertarget\hypertarget{sourcepointer-27}{}\fi
\noindent\textit{Source reference.} Original section 05N0: focused-seeds/AGP/sources/stacks-snapshot-20261011/tags.txt, source lines 7289--7289. Stable Stacks tag for an expressly allowed foundational or notational input; the counted construction remains the two concrete core obligations. The retained local core is not replaced by this pointer. See standard\_prerequisites.\par\medskip
\ifdefined\hypertarget\hypertarget{sourcepointer-28}{}\fi
\noindent\textit{Source reference.} Original lemma 0DKW: focused-seeds/AGP/sources/stacks-snapshot-20261011/tags.txt, source lines 17026--17026. Stable Stacks tag for an expressly allowed foundational or notational input; the counted construction remains the two concrete core obligations. The retained local core is not replaced by this pointer. See standard\_prerequisites.\par\medskip
\ifdefined\hypertarget\hypertarget{sourcepointer-29}{}\fi
\noindent\textit{Source reference.} Original lemma 05F4: focused-seeds/AGP/sources/stacks-snapshot-20261011/tags.txt, source lines 7013--7013. Stable Stacks tag for an expressly allowed foundational or notational input; the counted construction remains the two concrete core obligations. The retained local core is not replaced by this pointer. See standard\_prerequisites.\par\medskip
\ifdefined\hypertarget\hypertarget{sourcepointer-30}{}\fi
\noindent\textit{Source reference.} Original lemma 07RT: focused-seeds/AGP/sources/stacks-snapshot-20261011/tags.txt, source lines 9878--9878. Stable Stacks tag for an expressly allowed foundational or notational input; the counted construction remains the two concrete core obligations. The retained local core is not replaced by this pointer. See standard\_prerequisites.\par\medskip
\ifdefined\hypertarget\hypertarget{sourcepointer-31}{}\fi
\noindent\textit{Source reference.} Original lemma 07SY: focused-seeds/AGP/sources/stacks-snapshot-20261011/tags.txt, source lines 9918--9918. Stable Stacks tag for an expressly allowed foundational or notational input; the counted construction remains the two concrete core obligations. The retained local core is not replaced by this pointer. See standard\_prerequisites.\par\medskip
\ifdefined\hypertarget\hypertarget{sourcepointer-32}{}\fi
\noindent\textit{Source reference.} Original lemma 07WN: focused-seeds/AGP/sources/stacks-snapshot-20261011/tags.txt, source lines 10048--10048. Stable Stacks tag for an expressly allowed foundational or notational input; the counted construction remains the two concrete core obligations. The retained local core is not replaced by this pointer. See standard\_prerequisites.\par\medskip
\ifdefined\hypertarget\hypertarget{sourcepointer-33}{}\fi
\noindent\textit{Source reference.} Original lemma 07W3: focused-seeds/AGP/sources/stacks-snapshot-20261011/tags.txt, source lines 10028--10028. Stable Stacks tag for an expressly allowed foundational or notational input; the counted construction remains the two concrete core obligations. The retained local core is not replaced by this pointer. See standard\_prerequisites.\par\medskip
\ifdefined\hypertarget\hypertarget{sourcepointer-34}{}\fi
\noindent\textit{Source reference.} Original lemma 0DL7: focused-seeds/AGP/sources/stacks-snapshot-20261011/tags.txt, source lines 17037--17037. Stable Stacks tag for an expressly allowed foundational or notational input; the counted construction remains the two concrete core obligations. The retained local core is not replaced by this pointer. See standard\_prerequisites.\par\medskip
\ifdefined\hypertarget\hypertarget{sourcepointer-35}{}\fi
\noindent\textit{Source reference.} Original lemma 0DL4: focused-seeds/AGP/sources/stacks-snapshot-20261011/tags.txt, source lines 17034--17034. Stable Stacks tag for an expressly allowed foundational or notational input; the counted construction remains the two concrete core obligations. The retained local core is not replaced by this pointer. See standard\_prerequisites.\par\medskip
\ifdefined\hypertarget\hypertarget{sourcepointer-36}{}\fi
\noindent\textit{Source reference.} Original section 07WY: focused-seeds/AGP/sources/stacks-snapshot-20261011/tags.txt, source lines 10058--10058. Stable Stacks tag for an expressly allowed foundational or notational input; the counted construction remains the two concrete core obligations. The retained local core is not replaced by this pointer. See standard\_prerequisites.\par\medskip
\ifdefined\hypertarget\hypertarget{sourcepointer-37}{}\fi
\noindent\textit{Source reference.} Original section 0CXN: focused-seeds/AGP/sources/stacks-snapshot-20261011/tags.txt, source lines 16213--16213. Stable Stacks tag for an expressly allowed foundational or notational input; the counted construction remains the two concrete core obligations. The retained local core is not replaced by this pointer. See standard\_prerequisites.\par\medskip
\ifdefined\hypertarget\hypertarget{sourcepointer-38}{}\fi
\noindent\textit{Source reference.} Original lemma 07WU: focused-seeds/AGP/sources/stacks-snapshot-20261011/tags.txt, source lines 10054--10054. Stable Stacks tag for an expressly allowed foundational or notational input; the counted construction remains the two concrete core obligations. The retained local core is not replaced by this pointer. See standard\_prerequisites.\par\medskip
\ifdefined\hypertarget\hypertarget{sourcepointer-39}{}\fi
\noindent\textit{Source reference.} Original lemma 07X2: focused-seeds/AGP/sources/stacks-snapshot-20261011/tags.txt, source lines 10062--10062. Stable Stacks tag for an expressly allowed foundational or notational input; the counted construction remains the two concrete core obligations. The retained local core is not replaced by this pointer. See standard\_prerequisites.\par\medskip
\ifdefined\hypertarget\hypertarget{sourcepointer-40}{}\fi
\noindent\textit{Source reference.} Original lemma 07X1: focused-seeds/AGP/sources/stacks-snapshot-20261011/tags.txt, source lines 10061--10061. Stable Stacks tag for an expressly allowed foundational or notational input; the counted construction remains the two concrete core obligations. The retained local core is not replaced by this pointer. See standard\_prerequisites.\par\medskip
\ifdefined\hypertarget\hypertarget{sourcepointer-41}{}\fi
\noindent\textit{Source reference.} Original section 05ZJ: focused-seeds/AGP/sources/stacks-snapshot-20261011/tags.txt, source lines 7693--7693. Stable Stacks tag for an expressly allowed foundational or notational input; the counted construction remains the two concrete core obligations. The retained local core is not replaced by this pointer. See standard\_prerequisites.\par\medskip
\ifdefined\hypertarget\hypertarget{sourcepointer-42}{}\fi
\noindent\textit{Source reference.} Original lemma 0DIZ: focused-seeds/AGP/sources/stacks-snapshot-20261011/tags.txt, source lines 16959--16959. Stable Stacks tag for an expressly allowed foundational or notational input; the counted construction remains the two concrete core obligations. The retained local core is not replaced by this pointer. See standard\_prerequisites.\par\medskip
\ifdefined\hypertarget\hypertarget{sourcepointer-43}{}\fi
\noindent\textit{Source reference.} Original lemma 08IK: focused-seeds/AGP/sources/stacks-snapshot-20261011/tags.txt, source lines 10818--10818. Stable Stacks tag for an expressly allowed foundational or notational input; the counted construction remains the two concrete core obligations. The retained local core is not replaced by this pointer. See standard\_prerequisites.\par\medskip
\ifdefined\hypertarget\hypertarget{sourcepointer-44}{}\fi
\noindent\textit{Source reference.} Original lemma 0D0T: focused-seeds/AGP/sources/stacks-snapshot-20261011/tags.txt, source lines 16323--16323. Stable Stacks tag for an expressly allowed foundational or notational input; the counted construction remains the two concrete core obligations. The retained local core is not replaced by this pointer. See standard\_prerequisites.\par\medskip
\ifdefined\hypertarget\hypertarget{sourcepointer-45}{}\fi
\noindent\textit{Source reference.} Original section 0CXT: focused-seeds/AGP/sources/stacks-snapshot-20261011/tags.txt, source lines 16218--16218. Stable Stacks tag for an expressly allowed foundational or notational input; the counted construction remains the two concrete core obligations. The retained local core is not replaced by this pointer. See standard\_prerequisites.\par\medskip
\ifdefined\hypertarget\hypertarget{sourcepointer-46}{}\fi
\noindent\textit{Source reference.} Original theorem 0DIQ: focused-seeds/AGP/sources/stacks-snapshot-20261011/tags.txt, source lines 16950--16950. Stable Stacks tag for an expressly allowed foundational or notational input; the counted construction remains the two concrete core obligations. The retained local core is not replaced by this pointer. See standard\_prerequisites.\par\medskip
\ifdefined\hypertarget\hypertarget{sourcepointer-47}{}\fi
\noindent\textit{Source reference.} Original lemma 0CT7: focused-seeds/AGP/sources/stacks-snapshot-20261011/tags.txt, source lines 16057--16057. Stable Stacks tag for an expressly allowed foundational or notational input; the counted construction remains the two concrete core obligations. The retained local core is not replaced by this pointer. See standard\_prerequisites.\par\medskip
\ifdefined\hypertarget\hypertarget{sourcepointer-48}{}\fi
\noindent\textit{Source reference.} Original lemma 0DL5: focused-seeds/AGP/sources/stacks-snapshot-20261011/tags.txt, source lines 17035--17035. Stable Stacks tag for an expressly allowed foundational or notational input; the counted construction remains the two concrete core obligations. The retained local core is not replaced by this pointer. See standard\_prerequisites.\par\medskip
\ifdefined\hypertarget\hypertarget{sourcepointer-49}{}\fi
\noindent\textit{Source reference.} Original section 03YJ: focused-seeds/AGP/sources/stacks-snapshot-20261011/tags.txt, source lines 5186--5186. Stable Stacks tag for an expressly allowed foundational or notational input; the counted construction remains the two concrete core obligations. The retained local core is not replaced by this pointer. See standard\_prerequisites.\par\medskip
\ifdefined\hypertarget\hypertarget{sourcepointer-50}{}\fi
\noindent\textit{Source reference.} Original section 04X4: focused-seeds/AGP/sources/stacks-snapshot-20261011/tags.txt, source lines 6382--6382. Stable Stacks tag for an expressly allowed foundational or notational input; the counted construction remains the two concrete core obligations. The retained local core is not replaced by this pointer. See standard\_prerequisites.\par\medskip
\ifdefined\hypertarget\hypertarget{sourcepointer-51}{}\fi
\noindent\textit{Source reference.} Original lemma 07SK: focused-seeds/AGP/sources/stacks-snapshot-20261011/tags.txt, source lines 9905--9905. Stable Stacks tag for an expressly allowed foundational or notational input; the counted construction remains the two concrete core obligations. The retained local core is not replaced by this pointer. See standard\_prerequisites.\par\medskip
\ifdefined\hypertarget\hypertarget{sourcepointer-52}{}\fi
\noindent\textit{Source reference.} Original lemma 08K0: focused-seeds/AGP/sources/stacks-snapshot-20261011/tags.txt, source lines 10868--10868. Stable Stacks tag for an expressly allowed foundational or notational input; the counted construction remains the two concrete core obligations. The retained local core is not replaced by this pointer. See standard\_prerequisites.\par\medskip
\ifdefined\hypertarget\hypertarget{sourcepointer-53}{}\fi
\noindent\textit{Source reference.} Original lemma 08K1: focused-seeds/AGP/sources/stacks-snapshot-20261011/tags.txt, source lines 10869--10869. Stable Stacks tag for an expressly allowed foundational or notational input; the counted construction remains the two concrete core obligations. The retained local core is not replaced by this pointer. See standard\_prerequisites.\par\medskip
\ifdefined\hypertarget\hypertarget{sourcepointer-54}{}\fi
\noindent\textit{Source reference.} Original lemma 07Y4: focused-seeds/AGP/sources/stacks-snapshot-20261011/tags.txt, source lines 10099--10099. Stable Stacks tag for an expressly allowed foundational or notational input; the counted construction remains the two concrete core obligations. The retained local core is not replaced by this pointer. See standard\_prerequisites.\par\medskip
\ifdefined\hypertarget\hypertarget{sourcepointer-55}{}\fi
\noindent\textit{Source reference.} Original proposition 07PX: focused-seeds/AGP/sources/stacks-snapshot-20261011/tags.txt, source lines 9812--9812. Stable Stacks tag for an expressly allowed foundational or notational input; the counted construction remains the two concrete core obligations. The retained local core is not replaced by this pointer. See standard\_prerequisites.\par\medskip
\ifdefined\hypertarget\hypertarget{sourcepointer-56}{}\fi
\noindent\textit{Source reference.} Original section 07XJ: focused-seeds/AGP/sources/stacks-snapshot-20261011/tags.txt, source lines 10079--10079. Stable Stacks tag for an expressly allowed foundational or notational input; the counted construction remains the two concrete core obligations. The retained local core is not replaced by this pointer. See standard\_prerequisites.\par\medskip
\ifdefined\hypertarget\hypertarget{sourcepointer-57}{}\fi
\noindent\textit{Source reference.} Original lemma 0CXU: focused-seeds/AGP/sources/stacks-snapshot-20261011/tags.txt, source lines 16219--16219. Stable Stacks tag for an expressly allowed foundational or notational input; the counted construction remains the two concrete core obligations. The retained local core is not replaced by this pointer. See standard\_prerequisites.\par\medskip
\begin{theorem}[Algebraicity of moduli of complexes on a proper morphism]
\label{theorem-complexes-algebraic}
\begin{reference}
\cite{lieblich-complexes}
\end{reference}
Let $S$ be a scheme. Let $f : X \to B$ be morphism of algebraic spaces
over $S$. Assume that $f$ is proper, flat, and of finite presentation.
Then $\Complexesstack_{X/B}$ is an algebraic stack over $S$.
\end{theorem}

\section{Moduli of complexes on a proper morphism}
\label{section-moduli-complexes}

\noindent
The title and the material of this section are taken from
\cite{lieblich-complexes}. Let $S$ be a scheme and let
$f : X \to B$ be a proper, flat, finitely presented morphism
of algebraic spaces. We will prove that there is an
algebraic stack
$$
\Complexesstack_{X/B}
$$
parametrizing ``families'' of objects of $D^b_{\textit{Coh}}$
of the fibres with vanishing negative self-exts. More precisely
a family is given by a relatively perfect object of the derived
category of the total space; this somewhat technical notion
is studied in
More on Morphisms of Spaces, Section
\ref{spaces-more-morphisms-section-relatively-perfect}.

\medskip\noindent
Already if $X$ is a proper algebraic
space over a field $k$ we obtain a very interesting algebraic stack.
Namely, there is an embedding
$$
\Cohstack_{X/k} \longrightarrow \Complexesstack_{X/k}
$$
since for any $\mathcal{O}$-module $\mathcal{F}$ (on any ringed topos)
we have $\Ext^i_\mathcal{O}(\mathcal{F}, \mathcal{F}) = 0$ for $i < 0$.
Although this certainly shows our stack is nonempty, the
true motivation for the study of $\Complexesstack_{X/k}$
is that there are often objects of the derived category
$D^b_{\textit{Coh}}(\mathcal{O}_X)$ with vanishing negative self-exts
and nonvanishing cohomology sheaves in more than one degree.
For example, $X$ could be derived equivalent to another
proper algebraic space $Y$ over $k$, i.e., we have a
$k$-linear equivalence
$$
F : D^b_{\textit{Coh}}(\mathcal{O}_Y)
\longrightarrow
D^b_{\textit{Coh}}(\mathcal{O}_X)
$$
There are cases where this happens and $F$ is not given by
an isomorphism between $X$ and $Y$; for example in the case
of an abelian variety and its dual. In this situation $F$ induces
an isomorphism of algebraic stacks
$$
\Complexesstack_{Y/k}
\longrightarrow
\Complexesstack_{X/k}
$$
(insert future reference here) and in particular the stack of coherent sheaves
on $Y$ maps into the stack of complexes on $X$. Turning this around,
if we can understand well enough the geometry of
$\Complexesstack_{X/k}$, then we can try to use this to study all possible
derived equivalent $Y$.

\begin{lemma}
\label{lemma-complexes-open-neg-exts-vanishing}
Let $S$ be a scheme.
Let $f : X \to Y$ be a morphism of algebraic spaces over $S$.
Assume $f$ is proper, flat, and of finite presentation.
Let $K, E \in D(\mathcal{O}_X)$. Assume $K$ is pseudo-coherent
and $E$ is $Y$-perfect (More on Morphisms of Spaces, Definition
\ref{spaces-more-morphisms-definition-relatively-perfect}).
For a field $k$ and a morphism $y : \Spec(k) \to Y$ denote $K_y$, $E_y$
the pullback to the fibre $X_y$.
\begin{enumerate}
\item There is an open $W \subset Y$ characterized by the property
$$
y \in |W|
\Leftrightarrow
\Ext^i_{\mathcal{O}_{X_y}}(K_y, E_y) = 0
\text{ for }i < 0.
$$
\item For any morphism $V \to Y$ factoring through $W$ we have
$$
\Ext^i_{\mathcal{O}_{X_V}}(K_V, E_V) = 0
\quad\text{for}\quad i < 0
$$
where $X_V$ is the base change of $X$ and $K_V$ and $E_V$
are the derived pullbacks of $K$ and $E$ to $X_V$.
\item The functor $V \mapsto \Hom_{\mathcal{O}_{X_V}}(K_V, E_V)$
is a sheaf on $(\textit{Spaces}/W)_{fppf}$ representable by an
algebraic space affine and of finite presentation over $W$.
\end{enumerate}
\end{lemma}

\begin{proof}
For any morphism $V \to Y$ the complex $K_V$ is pseudo-coherent
(Cohomology on Sites, Lemma
\ref{sites-cohomology-lemma-pseudo-coherent-pullback})
and $E_V$ is $V$-perfect (More on Morphisms of Spaces, Lemma
\ref{spaces-more-morphisms-lemma-base-change-relatively-perfect}).
Another observation is that given $y : \Spec(k) \to Y$
and a field extension $k'/k$ with $y' : \Spec(k') \to Y$
the induced morphism, we have
$$
\Ext^i_{\mathcal{O}_{X_{y'}}}(K_{y'}, E_{y'}) =
\Ext^i_{\mathcal{O}_{X_y}}(K_y, E_y) \otimes_k k'
$$
by Derived Categories of Schemes, Lemma
\ref{perfect-lemma-affine-morphism-and-hom-out-of-perfect}.
Thus the vanishing in (1) is really a property of the induced
point $y \in |Y|$.
We will use these two observations without further mention in the proof.

\medskip\noindent
Assume first $Y$ is an affine scheme. Then we may apply
More on Morphisms of Spaces, Lemma
\ref{spaces-more-morphisms-lemma-compute-ext-rel-perfect}
and find a pseudo-coherent $L \in D(\mathcal{O}_Y)$ which
``universally computes'' $Rf_*R\SheafHom(K, E)$ in the sense
described in that lemma. Unwinding the definitions, we obtain
for a point $y \in Y$ the equality
$$
\Ext^i_{\kappa(y)}(L \otimes_{\mathcal{O}_Y}^\mathbf{L} \kappa(y),
\kappa(y)) = \Ext^i_{\mathcal{O}_{X_y}}(K_y, E_y)
$$
We conclude that
$$
H^i(L \otimes_{\mathcal{O}_Y}^\mathbf{L} \kappa(y)) = 0
\text{ for } i > 0 \Leftrightarrow
\Ext^i_{\mathcal{O}_{X_y}}(K_y, E_y) = 0 \text{ for }i < 0.
$$
By Derived Categories of Schemes, Lemma \ref{perfect-lemma-jump-loci}
the set $W$ of $y \in Y$ where this happens defines an open of $Y$.
This open $W$ then satisfies the requirement in (1) for all morphisms
from spectra of fields, by the ``universality'' of $L$.

\medskip\noindent
Let's go back to $Y$ a general algebraic space.
Choose an \'etale covering $\{V_i \to Y\}$ by affine schemes $V_i$.
Then we see that the subset $W \subset |Y|$ pulls back to the corresponding
subset $W_i \subset |V_i|$ for $X_{V_i}$, $K_{V_i}$, $E_{V_i}$.
By the previous paragraph we find that $W_i$ is open, hence $W$ is open.
This proves (1) in general. Moreover, parts (2) and (3) are entirely formulated
in terms of the category $\textit{Spaces}/W$ and the restrictions
$X_W$, $K_W$, $E_W$. This reduces us to the case $W = Y$.

\medskip\noindent
Assume $W = Y$. We claim that for any algebraic space $V$ over $Y$
we have $Rf_{V, *}R\SheafHom(K_V, E_V)$ has vanishing cohomology
sheaves in degrees $< 0$. This will prove (2) because
$$
\Ext^i_{\mathcal{O}_{X_V}}(K_V, E_V) =
H^i(X_V, R\SheafHom(K_V, E_V)) = 
H^i(V, Rf_{V, *}R\SheafHom(K_V, E_V))
$$
by Cohomology on Sites, Lemmas
\ref{sites-cohomology-lemma-section-RHom-over-U} and
\ref{sites-cohomology-lemma-Leray-unbounded}
and the vanishing of the cohomology sheaves implies the
cohomology group $H^i$ is zero for $i < 0$ by
Derived Categories, Lemma \ref{derived-lemma-negative-vanishing}.

\medskip\noindent
To prove the claim, we may work \'etale locally on $V$.
In particular, we may assume $Y$ is affine and $W = Y$.
Let $L \in D(\mathcal{O}_Y)$ be as in the second paragraph of the proof.
For an algebraic space $V$ over $Y$ denote $L_V$ the derived pullback of
$L$ to $V$. (An important feature we will use is that $L$ ``works'' for all
algebraic spaces $V$ over $Y$ and not just affine $V$.)
As $W = Y$ we have $H^i(L) = 0$ for $i > 0$
(use More on Algebra, Lemma
\ref{more-algebra-lemma-lift-pseudo-coherent-from-residue-field}
to go from fibres to stalks). Hence $H^i(L_V) = 0$ for $i > 0$.
The property defining $L$ is that
$$
Rf_{V, *}R\SheafHom(K_V, E_V) = R\SheafHom(L_V, \mathcal{O}_V)
$$
Since $L_V$ sits in degrees $\leq 0$, we conclude that
$R\SheafHom(L_V, \mathcal{O}_V)$ sits in degrees $\geq 0$
thereby proving the claim. This finishes the proof of (2).

\medskip\noindent
Assume $W = Y$ but make no assumptions on the algebraic space $Y$.
Since we have (2), we see from
Simplicial Spaces, Lemma \ref{spaces-simplicial-lemma-fppf-neg-ext-zero-hom}
that the functor $F$ given by $F(V) = \Hom_{\mathcal{O}_{X_V}}(K_V, E_V)$
is a sheaf\footnote{To check the sheaf property
for a covering $\{V_i \to V\}_{i \in I}$ first consider the
{\v C}ech fppf hypercovering $a : V_\bullet \to V$ with
$V_n = \coprod_{i_0 \ldots i_n} V_{i_0} \times_V \ldots \times_V V_{i_n}$
and then set $U_\bullet = V_\bullet \times_{a, V} X_V$. Then
$U_\bullet \to X_V$ is an fppf hypercovering to which we may
apply Simplicial Spaces, Lemma
\ref{spaces-simplicial-lemma-fppf-neg-ext-zero-hom}.}
on $(\textit{Spaces}/Y)_{fppf}$. Thus to prove that $F$
is an algebraic space and that $F \to Y$ is affine and of
finite presentation, we may work \'etale locally on $Y$; see
Bootstrap, Lemma \ref{bootstrap-lemma-locally-algebraic-space-finite-type}
and
Morphisms of Spaces, Lemmas \ref{spaces-morphisms-lemma-affine-local} and
\ref{spaces-morphisms-lemma-finite-presentation-local}. We conclude
that it suffices to prove $F$ is an affine algebraic space of
finite presentation over $Y$ when $Y$ is an affine scheme. In this
case we go back to our pseudo-coherent complex $L \in D(\mathcal{O}_Y)$.
Since $H^i(L) = 0$ for $i > 0$, we can represent $L$ by a complex
of the form
$$
\ldots \to \mathcal{O}_Y^{\oplus m_1} \to \mathcal{O}_Y^{\oplus m_0}
\to 0 \to \ldots
$$
with the last term in degree $0$, see More on Algebra, Lemma
\ref{more-algebra-lemma-pseudo-coherent}. Combining the two displayed formulas
earlier in the proof we find that
$$
F(V) =
\Ker(
\Hom_V(\mathcal{O}_V^{\oplus m_0}, \mathcal{O}_V)
\to 
\Hom_V(\mathcal{O}_V^{\oplus m_1}, \mathcal{O}_V)
)
$$
In other words, there is a fibre product diagram
$$
\xymatrix{
F \ar[d] \ar[r] & Y \ar[d]^0 \\
\mathbf{A}_Y^{m_0} \ar[r] & \mathbf{A}_Y^{m_1}
}
$$
which proves what we want.
\end{proof}

\begin{lemma}
\label{lemma-complexes}
Let $S$ be a scheme. Let $f : X \to Y$ be a morphism of algebraic
spaces over $S$. Assume $f$ is proper, flat, and of finite presentation.
Let $E \in D(\mathcal{O}_X)$.
Assume
\begin{enumerate}
\item $E$ is $S$-perfect (More on Morphisms of Spaces, Definition
\ref{spaces-more-morphisms-definition-relatively-perfect}), and
\item for every point $s \in S$ we have
$$
\Ext^i_{\mathcal{O}_{X_s}}(E_s, E_s) = 0
\quad\text{for}\quad i < 0
$$
where $E_s$ is the pullback to the fibre $X_s$.
\end{enumerate}
Then
\begin{enumerate}
\item[(a)] (1) and (2) are preserved by arbitrary base change $V \to Y$,
\item[(b)] $\Ext^i_{\mathcal{O}_{X_V}}(E_V, E_V) = 0$ for $i < 0$
and all $V$ over $Y$,
\item[(c)] $V \mapsto \Hom_{\mathcal{O}_{X_V}}(E_V, E_V)$ is representable
by an algebraic space affine and of finite presentation over $Y$.
\end{enumerate}
Here $X_V$ is the base change of $X$ and $E_V$ is the derived pullback
of $E$ to $X_V$.
\end{lemma}

\begin{proof}
Immediate consequence of Lemma \ref{lemma-complexes-open-neg-exts-vanishing}.
\end{proof}

\begin{situation}
\label{situation-complexes}
Let $S$ be a scheme. Let $f : X \to B$ be a morphism of algebraic spaces
over $S$. Assume $f$ is proper, flat, and of finite presentation.
We denote $\Complexesstack_{X/B}$ the category whose objects are
triples $(T, g, E)$ where
\begin{enumerate}
\item $T$ is a scheme over $S$,
\item $g : T \to B$ is a morphism over $S$, and setting
$X_T = T \times_{g, B} X$
\item $E$ is an object of $D(\mathcal{O}_{X_T})$ satisfying
conditions (1) and (2) of Lemma \ref{lemma-complexes}.
\end{enumerate}
A morphism $(T, g, E) \to (T', g', E')$
is given by a pair $(h, \varphi)$ where
\begin{enumerate}
\item $h : T \to T'$ is a morphism of schemes over $B$
(i.e., $g' \circ h = g$), and
\item $\varphi : L(h')^*E' \to E$ is an
isomorphism of $D(\mathcal{O}_{X_T})$ where $h' : X_T \to X_{T'}$
is the base change of $h$.
\end{enumerate}
\end{situation}

\noindent
Thus $\Complexesstack_{X/B}$ is a category and the rule
$$
p : \Complexesstack_{X/B} \longrightarrow (\Sch/S)_{fppf},
\quad
(T, g, E) \longmapsto T
$$
is a functor. For a scheme $T$ over $S$ we denote $\Complexesstack_{X/B, T}$
the fibre category of $p$ over $T$. These fibre categories are groupoids.

\begin{lemma}
\label{lemma-complexes-fibred-in-groupoids}
In Situation \ref{situation-complexes} the functor
$p : \Complexesstack_{X/B} \longrightarrow (\Sch/S)_{fppf}$
is fibred in groupoids.
\end{lemma}

\begin{proof}
We show that $p$ is fibred in groupoids by checking conditions
(1) and (2) of Categories, Definition
\ref{categories-definition-fibred-groupoids}.
Given an object $(T', g', E')$
of $\Complexesstack_{X/B}$ and a morphism $h : T \to T'$ of
schemes over $S$ we can set $g = h \circ g'$ and
$E = L(h')^*E'$ where $h' : X_T \to X_{T'}$
is the base change of $h$. Then it is clear that we obtain
a morphism $(T, g, E) \to (T', g', E')$
of $\Complexesstack_{X/B}$ lying over $h$. This proves (1).
For (2) suppose we are given morphisms
$$
(h_1, \varphi_1) : (T_1, g_1, E_1) \to (T, g, E)
\quad\text{and}\quad
(h_2, \varphi_2) : (T_2, g_2, E_2) \to (T, g, E)
$$
of $\Complexesstack_{X/B}$ and a morphism $h : T_1 \to T_2$ such that
$h_2 \circ h = h_1$. Then we can let $\varphi$ be the composition
$$
L(h')^*E_2
\xrightarrow{L(h')^*\varphi_2^{-1}}
L(h')^*L(h_2)^*E = L(h_1)^*E
\xrightarrow{\varphi_1}
E_1
$$
to obtain the morphism
$(h, \varphi) : (T_1, g_1, E_1) \to (T_2, g_2, E_2)$
that witnesses the truth of condition (2).
\end{proof}

\begin{lemma}
\label{lemma-complexes-diagonal}
In Situation \ref{situation-complexes}. Denote
$\mathcal{X} = \Complexesstack_{X/B}$. Then
$\Delta : \mathcal{X} \to \mathcal{X} \times \mathcal{X}$ is
representable by algebraic spaces.
\end{lemma}

\begin{proof}
Consider two objects $x = (T, g, E)$ and $y = (T, g', E')$
of $\mathcal{X}$ over a scheme $T$. We have to show that
$\mathit{Isom}_\mathcal{X}(x, y)$ is an algebraic space over $T$, see
Algebraic Stacks, Lemma \ref{algebraic-lemma-representable-diagonal}.
If for $h : T' \to T$ the restrictions $x|_{T'}$ and $y|_{T'}$ are isomorphic
in the fibre category $\mathcal{X}_{T'}$, then $g \circ h = g' \circ h$.
Hence there is a transformation of presheaves
$$
\mathit{Isom}_\mathcal{X}(x, y) \longrightarrow \text{Equalizer}(g, g')
$$
Since the diagonal of $B$ is representable (by schemes) this equalizer is
a scheme. Thus we may replace $T$ by this equalizer and
$E$ and $E'$ by their pullbacks. Thus we may assume $g = g'$.

\medskip\noindent
Assume $g = g'$. After replacing $B$ by $T$ and $X$ by $X_T$ we arrive
at the following problem. Given $E, E' \in D(\mathcal{O}_X)$
satisfying conditions (1), (2) of Lemma \ref{lemma-complexes}
we have to show that $\mathit{Isom}(E, E')$ is an algebraic space.
Here $\mathit{Isom}(E, E')$ is the functor
$$
(\Sch/B)^{opp} \to \textit{Sets},\quad
T \mapsto \{\varphi : E_T \to E'_T
\text{ isomorphism in }D(\mathcal{O}_{X_T})\}
$$
where $E_T$ and $E'_T$ are the derived pullbacks of $E$ and $E'$
to $X_T$. Now, let $W \subset B$, resp.\ $W' \subset B$ be the
open subspace of $B$ associated
to $E, E'$, resp.\ to $E', E$ by
Lemma \ref{lemma-complexes-open-neg-exts-vanishing}.
Clearly, if there exists an isomorphism $E_T \to E'_T$ as in
the definition of $\mathit{Isom}(E, E')$, then we see that
$T \to B$ factors into both $W$ and $W'$ (because we have
condition (1) for $E$ and $E'$ and we'll obviously have
$E_t \cong E'_t$ so no nonzero maps $E_t[i] \to E_t$
or $E'_t[i] \to E_t$ over the fibre $X_t$ for $i > 0$.
Thus we may replace $B$ by the open $W \cap W'$.
In this case the functor $H = \SheafHom(E, E')$
$$
(\Sch/B)^{opp} \to \textit{Sets},\quad T \mapsto
\Hom_{\mathcal{O}_{X_T}}(E_T, E'_T)
$$
is an algebraic space affine and of finite presentation over $B$ by
Lemma \ref{lemma-complexes-open-neg-exts-vanishing}.
The same is true for
$H' = \SheafHom(E', E)$,
$I = \SheafHom(E, E)$, and
$I' = \SheafHom(E', E')$.
Therefore we can repeat the argument of the proof of
Proposition \ref{proposition-isom} to see that
$$
\mathit{Isom}(E, E') = (H' \times_B H) \times_{c, I \times_B I', \sigma} B
$$
for some morphisms $c$ and $\sigma$. Thus
$\mathit{Isom}(E, E')$ is an algebraic space.
\end{proof}

\begin{lemma}
\label{lemma-complexes-stack}
In Situation \ref{situation-complexes} the functor
$p : \Complexesstack_{X/B} \longrightarrow (\Sch/S)_{fppf}$
is a stack in groupoids.
\end{lemma}

\begin{proof}
To prove that $\Complexesstack_{X/B}$ is a stack in groupoids,
we have to show that the presheaves $\mathit{Isom}$ are sheaves
and that descent data are effective. The statement on
$\mathit{Isom}$ follows from Lemma \ref{lemma-complexes-diagonal}, see
Algebraic Stacks, Lemma \ref{algebraic-lemma-representable-diagonal}.
Let us prove the statement on descent data.

\medskip\noindent
Suppose that $\{a_i : T_i \to T\}$ is an fppf covering of schemes over $S$.
Let $(\xi_i, \varphi_{ij})$ be a descent datum for $\{T_i \to T\}$
with values in $\Complexesstack_{X/B}$.
For each $i$ we can write $\xi_i = (T_i, g_i, E_i)$.
Denote $\text{pr}_0 : T_i \times_T T_j \to T_i$ and
$\text{pr}_1 : T_i \times_T T_j \to T_j$ the projections.
The condition that $\xi_i|_{T_i \times_T T_j} \cong \xi_j|_{T_i \times_T T_j}$
implies in particular that $g_i \circ \text{pr}_0 = g_j \circ \text{pr}_1$.
Thus there exists a unique morphism $g : T \to B$ such that
$g_i = g \circ a_i$, see
Descent on Spaces, Lemma
\ref{spaces-descent-lemma-fpqc-universal-effective-epimorphisms}.
Denote $X_T = T \times_{g, B} X$. Set
$X_i = X_{T_i} = T_i \times_{g_i, B} X = T_i \times_{a_i, T} X_T$
and
$$
X_{ij} = X_{T_i} \times_{X_T} X_{T_j} = X_i \times_{X_T} X_j
$$
with projections $\text{pr}_i$ and $\text{pr}_j$ to $X_i$ and $X_j$.
Observe that the pullback of $(T_i, g_i, E_i)$
by $\text{pr}_0 : T_i \times_T T_j \to T_i$ is given by
$(T_i \times_T T_j, g_i \circ \text{pr}_0, L\text{pr}_i^*E_i)$.
Hence a descent datum for $\{T_i \to T\}$ in $\Complexesstack_{X/B}$
is given by the objects $(T_i, g \circ a_i, E_i)$
and for each pair $i, j$ an isomorphism in
$D\mathcal{O}_{X_{ij}})$
$$
\varphi_{ij} :
L\text{pr}_i^*E_i \longrightarrow L\text{pr}_j^*E_j
$$
satisfying the cocycle condition over the pullback of $X$ to
$T_i \times_T T_j \times_T T_k$.
Using the vanishing of negative Exts provided by (b) of
Lemma \ref{lemma-complexes}, we may apply
Simplicial Spaces, Lemma \ref{spaces-simplicial-lemma-fppf-glue-neg-ext-zero}
to obtain descent\footnote{To check this, first consider the
{\v C}ech fppf hypercovering $a : T_\bullet \to T$ with
$T_n = \coprod_{i_0 \ldots i_n} T_{i_0} \times_T \ldots \times_T T_{i_n}$
and then set $U_\bullet = T_\bullet \times_{a, T} X_T$. Then
$U_\bullet \to X_T$ is an fppf hypercovering to which we may
apply Simplicial Spaces, Lemma
\ref{spaces-simplicial-lemma-fppf-glue-neg-ext-zero}.}
for these complexes. In other words, we find there exists an object
$E$ in $D_\QCoh(\mathcal{O}_{X_T})$ restricting to $E_i$ on $X_{T_i}$
compatible with $\varphi_{ij}$. Recall that being
$T$-perfect signifies being pseudo-coherent and having
locally finite tor dimension over $f^{-1}\mathcal{O}_T$.
Thus $E$ is $T$-perfect by an application of
More on Morphisms of Spaces, Lemmas
\ref{spaces-more-morphisms-lemma-pseudo-coherent-descends-fpqc} and
\ref{spaces-more-morphisms-lemma-tor-amplitude-descends-fppf}.
Finally, we have to check condition (2) from
Lemma \ref{lemma-complexes} for $E$.
This immediately follows from the description of the open $W$
in Lemma \ref{lemma-complexes-open-neg-exts-vanishing}
and the fact that (2) holds for $E_i$ on $X_{T_i}/T_i$.
\end{proof}

\begin{remark}
\label{remark-complexes-base-change}
In Situation \ref{situation-complexes} the rule
$(T, g, E) \mapsto (T, g)$ defines a $1$-morphism
$$
\Complexesstack_{X/B} \longrightarrow \mathcal{S}_B
$$
of stacks in groupoids
(see Lemma \ref{lemma-complexes-stack},
Algebraic Stacks, Section \ref{algebraic-section-split}, and
Examples of Stacks, Section
\ref{examples-stacks-section-stack-associated-to-sheaf}).
Let $B' \to B$ be a morphism of
algebraic spaces over $S$. Let $\mathcal{S}_{B'} \to \mathcal{S}_B$
be the associated $1$-morphism of stacks fibred in sets.
Set $X' = X \times_B B'$.
We obtain a stack in groupoids
$\Complexesstack_{X'/B'} \to (\Sch/S)_{fppf}$
associated to the base change $f' : X' \to B'$. In this situation
the diagram
$$
\vcenter{
\xymatrix{
\Complexesstack_{X'/B'} \ar[r] \ar[d] &
\Complexesstack_{X/B} \ar[d] \\
\mathcal{S}_{B'} \ar[r] & \mathcal{S}_B
}
}
\quad
\begin{matrix}
\text{or in} \\
\text{another} \\
\text{notation}
\end{matrix}
\quad
\vcenter{
\xymatrix{
\Complexesstack_{X'/B'} \ar[r] \ar[d] &
\Complexesstack_{X/B} \ar[d] \\
\Sch/B' \ar[r] & \Sch/B
}
}
$$
is $2$-fibre product square. This trivial remark
will occasionally be useful to change the base algebraic space.
\end{remark}

\begin{lemma}
\label{lemma-complexes-limits}
In Situation \ref{situation-complexes} assume that $B \to S$
is locally of finite presentation. Then
$p : \Complexesstack_{X/B} \to (\Sch/S)_{fppf}$ is limit preserving
(Artin's Axioms, Definition \ref{artin-definition-limit-preserving}).
\end{lemma}

\begin{proof}
Write $B(T)$ for the discrete category whose
objects are the $S$-morphisms $T \to B$. Let $T = \lim T_i$ be a filtered
limit of affine schemes over $S$. Assigning to an object
$(T, h, E)$ of $\Complexesstack_{X/B, T}$ the object $h$
of $B(T)$ gives us a commutative diagram of fibre categories
$$
\xymatrix{
\colim \Complexesstack_{X/B, T_i} \ar[r] \ar[d] &
\Complexesstack_{X/B, T} \ar[d] \\
\colim B(T_i) \ar[r] & B(T)
}
$$
We have to show the top horizontal arrow is an equivalence. Since
we have assume that $B$ is locally of finite presentation over $S$
we see from
Limits of Spaces, Remark \ref{spaces-limits-remark-limit-preserving}
that the bottom horizontal arrow is an equivalence. This means that
we may assume $T = \lim T_i$ be a filtered limit of affine schemes over
$B$. Denote $g_i : T_i \to B$ and $g : T \to B$ the corresponding
morphisms. Set $X_i = T_i \times_{g_i, B} X$ and $X_T = T \times_{g, B} X$.
Observe that $X_T = \colim X_i$.
By More on Morphisms of Spaces, Lemma
\ref{spaces-more-morphisms-lemma-descend-relatively-perfect}
the category of $T$-perfect objects of $D(\mathcal{O}_{X_T})$
is the colimit of the categories of $T_i$-perfect objects
of $D(\mathcal{O}_{X_{T_i}})$.
Thus all we have to prove is that given an $T_i$-perfect object
$E_i$ of $D(\mathcal{O}_{X_{T_i}})$ such that
the derived pullback $E$ of $E_i$ to $X_T$ satisfies
condition (2) of Lemma \ref{lemma-complexes},
then after increasing $i$ we have that
$E_i$ satisfies
condition (2) of Lemma \ref{lemma-complexes}.
Let $W \subset |T_i|$ be the open constructed
in Lemma \ref{lemma-complexes-open-neg-exts-vanishing}
for $E_i$ and $E_i$. By assumption on $E$ we find
that $T \to T_i$ factors through $T$.
Hence there is an $i' \geq i$
such that $T_{i'} \to T_i$ factors through $W$, see
Limits, Lemma \ref{limits-lemma-limit-contained-in-constructible}
Then $i'$ works by construction of $W$.
\end{proof}

\begin{lemma}
\label{lemma-complexes-RS-star}
In Situation \ref{situation-complexes}. Let
$$
\xymatrix{
Z \ar[r] \ar[d] & Z' \ar[d] \\
Y \ar[r] & Y'
}
$$
be a pushout in the category of schemes over $S$ where
$Z \to Z'$ is a finite order thickening and $Z \to Y$ is affine, see
More on Morphisms, Lemma \ref{more-morphisms-lemma-pushout-along-thickening}.
Then the functor on fibre categories
$$
\Complexesstack_{X/B, Y'}
\longrightarrow
\Complexesstack_{X/B, Y}
\times_{\Complexesstack_{X/B, Z}}
\Complexesstack_{X/B, Z'}
$$
is an equivalence.
\end{lemma}

\begin{proof}
Observe that the corresponding map
$$
B(Y') \longrightarrow B(Y) \times_{B(Z)} B(Z')
$$
is a bijection, see Pushouts of Spaces, Lemma
\ref{spaces-pushouts-lemma-pushout-along-thickening-schemes}.
Thus using the commutative diagram
$$
\xymatrix{
\Complexesstack_{X/B, Y'} \ar[r] \ar[d] &
\Complexesstack_{X/B, Y}
\times_{\Complexesstack_{X/B, Z}}
\Complexesstack_{X/B, Z'}
\ar[d] \\
B(Y') \ar[r] & B(Y) \times_{B(Z)} B(Z')
}
$$
we see that we may assume that $Y'$ is a scheme over $B'$. By
Remark \ref{remark-complexes-base-change}
we may replace $B$ by $Y'$ and $X$ by $X \times_B Y'$.
Thus we may assume $B = Y'$.

\medskip\noindent
Assume $B = Y'$. We first prove fully faithfulness of our functor.
To do this, let $\xi_1, \xi_2$ be two objects of $\Complexesstack_{X/B}$
over $Y'$. Then we have to show that
$$
\mathit{Isom}(\xi_1, \xi_2)(Y') \longrightarrow
\mathit{Isom}(\xi_1, \xi_2)(Y)
\times_{\mathit{Isom}(\xi_1, \xi_2)(Z)}
\mathit{Isom}(\xi_1, \xi_2)(Z')
$$
is bijective. However, we already know that $\mathit{Isom}(\xi_1, \xi_2)$
is an algebraic space over $B = Y'$. Thus this bijectivity follows from
Artin's Axioms, Lemma \ref{artin-lemma-pushout} (or the aforementioned
Pushouts of Spaces, Lemma
\ref{spaces-pushouts-lemma-pushout-along-thickening-schemes}).

\medskip\noindent
Essential surjectivity. Let $(E_Y, E_{Z'}, \alpha)$ be a triple,
where $E_Y \in D(\mathcal{O}_Y)$ and $E_{Z'} \in D(\mathcal{O}_{X_{Z'}})$
are objects such that $(Y, Y \to B, E_Y)$ is an object of
$\Complexesstack_{X/B}$ over $Y$, such that
$(Z', Z' \to B, E_{Z'})$ is an object of $\Complexesstack_{X/B}$ over $Z'$,
and $\alpha : L(X_Z \to X_Y)^*E_Y \to L(X_Z \to X_{Z'})^*E_{Z'}$
is an isomorphism in $D(\mathcal{O}_{Z'})$.
That is to say
$$
((Y, Y \to B, E_Y), (Z', Z' \to B, E_{Z'}), \alpha)
$$
is an object of the target of the arrow of our lemma.
Observe that the diagram
$$
\xymatrix{
X_Z \ar[r] \ar[d] & X_{Z'} \ar[d] \\
X_Y \ar[r] & X_{Y'}
}
$$
is a pushout with $X_Z \to X_Y$ affine and $X_Z \to X_{Z'}$ a thickening
(see Pushouts of Spaces, Lemma
\ref{spaces-pushouts-lemma-equivalence-categories-spaces-pushout-flat}).
Hence by Pushouts of Spaces, Lemma
\ref{spaces-pushouts-lemma-pushout-along-thickening-derived}
we find an object $E_{Y'} \in D(\mathcal{O}_{X_{Y'}})$
together with isomorphisms
$L(X_Y \to X_{Y'})^*E_{Y'} \to E_Y$ and
$L(X_{Z'} \to X_{Y'})^*E_{Y'} \to E_Z$
compatible with $\alpha$. Clearly, if we show that
$E_{Y'}$ is $Y'$-perfect, then we are done, because property (2)
of Lemma \ref{lemma-complexes}
is a property on points (and $Y$ and $Y'$ have the same points).
This follows from More on Morphisms of Spaces, Lemma
\ref{spaces-more-morphisms-lemma-thickening-relatively-perfect}.
\end{proof}

\begin{lemma}
\label{lemma-complexes-tangent-space}
In Situation \ref{situation-complexes} assume that $S$ is a locally Noetherian
scheme and $B \to S$ is locally of finite presentation.
Let $k$ be a finite type field over $S$ and let
$x_0 = (\Spec(k), g_0, E_0)$
be an object of $\mathcal{X} = \Complexesstack_{X/B}$ over $k$.
Then the spaces $T\mathcal{F}_{\mathcal{X}, k, x_0}$ and
$\text{Inf}(\mathcal{F}_{\mathcal{X}, k, x_0})$
(Artin's Axioms, Section \ref{artin-section-tangent-spaces})
are finite dimensional.
\end{lemma}

\begin{proof}
Observe that by Lemma \ref{lemma-complexes-RS-star}
our stack in groupoids $\mathcal{X}$ satisfies property (RS*)
defined in Artin's Axioms, Section \ref{artin-section-RS-star}.
In particular $\mathcal{X}$ satisfies (RS).
Hence all associated predeformation
categories are deformation categories
(Artin's Axioms, Lemma \ref{artin-lemma-deformation-category})
and the statement makes sense.

\medskip\noindent
In this paragraph we show that we can reduce to the case $B = \Spec(k)$.
Set $X_0 = \Spec(k) \times_{g_0, B} X$
and denote $\mathcal{X}_0 = \Complexesstack_{X_0/k}$. In
Remark \ref{remark-complexes-base-change} we have seen that
$\mathcal{X}_0$ is the $2$-fibre product of $\mathcal{X}$ with
$\Spec(k)$ over $B$ as categories fibred in groupoids over
$(\Sch/S)_{fppf}$. Thus by
Artin's Axioms, Lemma \ref{artin-lemma-fibre-product-tangent-spaces}
we reduce to proving that $B$, $\Spec(k)$, and $\mathcal{X}_0$
have finite dimensional tangent spaces and infinitesimal automorphism
spaces. The tangent space of $B$ and $\Spec(k)$ are finite dimensional by
Artin's Axioms, Lemma \ref{artin-lemma-finite-dimension}
and of course these have vanishing $\text{Inf}$.
Thus it suffices to deal with $\mathcal{X}_0$.

\medskip\noindent
Let $k[\epsilon]$ be the dual numbers over $k$.
Let $\Spec(k[\epsilon]) \to B$ be the composition of $g_0 : \Spec(k) \to B$
and the morphism $\Spec(k[\epsilon]) \to \Spec(k)$
coming from the inclusion $k \to k[\epsilon]$.
Set $X_0 = \Spec(k) \times_B X$ and
$X_\epsilon = \Spec(k[\epsilon]) \times_B X$.
Observe that $X_\epsilon$ is a first order thickening of $X_0$
flat over the first order thickening $\Spec(k) \to \Spec(k[\epsilon])$.
Observe that $X_0$ and $X_\epsilon$ give rise to canonically equivalent
small \'etale topoi, see
More on Morphisms of Spaces, Section
\ref{spaces-more-morphisms-section-thickenings}.
By More on Morphisms of Spaces, Lemma
\ref{spaces-more-morphisms-lemma-thickening-relatively-perfect}
we see that $T\mathcal{F}_{\mathcal{X}_0, k, x_0}$ is the set of
isomorphism classes of lifts of $E_0$ to $X_\epsilon$ in the sense of
Deformation Theory, Lemma \ref{defos-lemma-first-order-defos-complex}.
We conclude that
$$
T\mathcal{F}_{\mathcal{X}_0, k, x_0} =
\Ext^1_{\mathcal{O}_{X_0}}(E_0, E_0)
$$
Here we have used the identification $\epsilon k[\epsilon] \cong k$
of $k[\epsilon]$-modules. Using
Deformation Theory, Lemma \ref{defos-lemma-first-order-defos-complex}
once more we see that there is a surjection
$$
\text{Inf}(\mathcal{F}_{\mathcal{X}, k, x_0})
\leftarrow
\Ext^0_{\mathcal{O}_{X_0}}(E_0, E_0)
$$
of $k$-vector spaces. As $E_0$ is pseudo-coherent it lies in
$D^-_{\textit{Coh}}(\mathcal{O}_{X_0})$ by
Derived Categories of Spaces, Lemma
\ref{spaces-perfect-lemma-identify-pseudo-coherent-noetherian}.
Since $E_0$ locally has finite tor dimension and $X_0$
is quasi-compact we see $E_0 \in D^b_{\textit{Coh}}(\mathcal{O}_{X_0})$.
Thus the $\Ext$s above are finite dimensional $k$-vector spaces
by Derived Categories of Spaces, Lemma
\ref{spaces-perfect-lemma-ext-finite}.
\end{proof}

\begin{lemma}
\label{lemma-complexes-strong-effectiveness}
In Situation \ref{situation-complexes} assume $B = S$ is locally Noetherian.
Then strong formal effectiveness in the sense of
Artin's Axioms, Remark \ref{artin-remark-strong-effectiveness}
holds for $p : \Complexesstack_{X/S} \to (\Sch/S)_{fppf}$.
\end{lemma}

\begin{proof}
Let $(R_n)$ be an inverse system of $S$-algebras with surjective transition
maps whose kernels are locally nilpotent. Set $R = \lim R_n$.
Let $(\xi_n)$ be a system of objects of $\Complexesstack_{X/B}$
lying over $(\Spec(R_n))$. We have to show $(\xi_n)$ is effective, i.e.,
there exists an object $\xi$ of $\Complexesstack_{X/B}$ lying over
$\Spec(R)$.

\medskip\noindent
Write $X_R = \Spec(R) \times_S X$ and $X_n = \Spec(R_n) \times_S X$.
Of course $X_n$ is the base change of $X_R$ by $R \to R_n$. Since $S = B$,
we see that $\xi_n$ corresponds simply to an $R_n$-perfect object
$E_n \in D(\mathcal{O}_{X_n})$ satisfying condition (2) of
Lemma \ref{lemma-complexes}. In particular $E_n$ is pseudo-coherent.
The isomorphisms $\xi_{n + 1}|_{\Spec(R_n)} \cong \xi_n$
correspond to isomorphisms $L(X_n \to X_{n + 1})^*E_{n + 1} \to E_n$.
Therefore by
Flatness on Spaces, Theorem \ref{spaces-flat-theorem-existence-derived}
we find a pseudo-coherent object $E$ of $D(\mathcal{O}_{X_R})$
with $E_n$ equal to the derived pullback of $E$ for all $n$
compatible with the transition isomorphisms.

\medskip\noindent
Observe that $(R, \Ker(R \to R_1))$ is a henselian pair, see
More on Algebra, Lemma \ref{more-algebra-lemma-limit-henselian}.
In particular, $\Ker(R \to R_1)$ is contained in the Jacobson radical of $R$.
Then we may apply More on Morphisms of Spaces, Lemma
\ref{spaces-more-morphisms-lemma-henselian-relatively-perfect}
to see that $E$ is $R$-perfect.

\medskip\noindent
Finally, we have to check condition (2) of Lemma \ref{lemma-complexes}.
By Lemma \ref{lemma-complexes-open-neg-exts-vanishing}
the set of points $t$ of $\Spec(R)$ where the negative self-exts of $E_t$
vanish is an open. Since this condition is true in $V(\Ker(R \to R_1))$
and since $\Ker(R \to R_1)$ is contained in the Jacobson radical of $R$
we conclude it holds for all points.
\end{proof}

\begin{proof}
Set $\mathcal{X} = \Complexesstack_{X/B}$. We have seen that $\mathcal{X}$
is a stack in groupoids over $(\Sch/S)_{fppf}$ with diagonal representable
by algebraic spaces
(Lemmas \ref{lemma-complexes-stack} and \ref{lemma-complexes-diagonal}).
Hence it suffices to find a scheme $W$ and a surjective and smooth
morphism $W \to \mathcal{X}$.

\medskip\noindent
Let $B'$ be a scheme and let $B' \to B$ be a surjective \'etale morphism.
Set $X' = B' \times_B X$ and denote $f' : X' \to B'$ the projection.
Then $\mathcal{X}' = \Complexesstack_{X'/B'}$ is equal to the $2$-fibre
product of $\mathcal{X}$ with the category fibred in sets
associated to $B'$ over the category fibred in sets associated to $B$
(Remark \ref{remark-complexes-base-change}). By the material in
Algebraic Stacks, Section \ref{algebraic-section-representable-properties}
the morphism $\mathcal{X}' \to \mathcal{X}$ is surjective and \'etale.
Hence it suffices to prove the result for $\mathcal{X}'$.
In other words, we may assume $B$ is a scheme.

\medskip\noindent
Assume $B$ is a scheme. In this case we may replace $S$ by $B$, see
Algebraic Stacks, Section \ref{algebraic-section-change-base-scheme}.
Thus we may assume $S = B$.

\medskip\noindent
Assume $S = B$. Choose an affine open covering $S = \bigcup U_i$.
Denote $\mathcal{X}_i$ the restriction of $\mathcal{X}$ to
$(\Sch/U_i)_{fppf}$. If we can find schemes $W_i$ over $U_i$ and
surjective smooth morphisms $W_i \to \mathcal{X}_i$, then we
set $W = \coprod W_i$ and we obtain a surjective smooth morphism
$W \to \mathcal{X}$. Thus we may assume $S = B$ is affine.

\medskip\noindent
Assume $S = B$ is affine, say $S = \Spec(\Lambda)$.
Write $\Lambda = \colim \Lambda_i$ as a filtered colimit with each $\Lambda_i$
of finite type over $\mathbf{Z}$. For some $i$ we can find
a morphism of algebraic spaces $X_i \to \Spec(\Lambda_i)$
which is proper, flat, of finite presentation and whose base change
to $\Lambda$ is $X$. See Limits of Spaces, Lemmas
\ref{spaces-limits-lemma-descend-finite-presentation},
\ref{spaces-limits-lemma-descend-flat}, and
\ref{spaces-limits-lemma-eventually-proper}.
If we show that $\Complexesstack_{X_i/\Spec(\Lambda_i)}$ is an
algebraic stack, then it follows by base change
(Remark \ref{remark-complexes-base-change} and
Algebraic Stacks, Section \ref{algebraic-section-change-base-scheme})
that $\mathcal{X}$ is an algebraic stack.
Thus we may assume that $\Lambda$ is a finite type $\mathbf{Z}$-algebra.

\medskip\noindent
Assume $S = B = \Spec(\Lambda)$ is affine of finite type over $\mathbf{Z}$.
In this case we will verify conditions (1), (2), (3), (4), and (5) of
Artin's Axioms, Lemma \ref{artin-lemma-diagonal-representable}
to conclude that $\mathcal{X}$ is an algebraic stack.
Note that $\Lambda$ is a G-ring, see
More on Algebra, Proposition \ref{more-algebra-proposition-ubiquity-G-ring}.
Hence all local rings of $S$ are G-rings. Thus (5) holds.
To check (2) we have to verify axioms [-1], [0], [1], [2], and [3]
of Artin's Axioms, Section \ref{artin-section-axioms}.
We omit the verification of [-1] and axioms
[0], [1], [2], [3] correspond respectively to
Lemmas \ref{lemma-complexes-stack},
\ref{lemma-complexes-limits},
\ref{lemma-complexes-RS-star},
\ref{lemma-complexes-tangent-space}.
Condition (3) follows from Lemma \ref{lemma-complexes-strong-effectiveness}.
Condition (1) is Lemma \ref{lemma-complexes-diagonal}.

\medskip\noindent
It remains to show condition (4) which is openness of versality.
To see this we will use
Artin's Axioms, Lemma \ref{artin-lemma-SGE-implies-openness-versality}.
We have already seen that $\mathcal{X}$ has diagonal
representable by algebraic spaces, has (RS*), and is limit preserving
(see lemmas used above).
Hence we only need to see that $\mathcal{X}$ satisfies the strong
formal effectiveness formulated in
Artin's Axioms, Lemma \ref{artin-lemma-SGE-implies-openness-versality}.
This follows from Lemma \ref{lemma-complexes-strong-effectiveness}
and the proof is complete.
\end{proof}

\bibliographystyle{plain}
\bibliography{my}
\end{document}
